% コンパイル（slide/20261107/ で実行。PDFは slide/out/20261107/ 、他は build/）:
%   ../lualatex-dated.sh -synctex=1 -interaction=nonstopmode -file-line-error -output-directory=out survey_kawamotostyle_20261107_handout.tex
\documentclass[aspectratio=169,10pt,t,handout]{beamer}
\usepackage[ja]{labslide}
\usepackage{multirow}
% 本文・見出し・脚注など、指定されている文字サイズを一律に 1pt 大きくする．
\makeatletter
\newdimen\labvtwo@size
\newskip\labvtwo@skip
\let\labvtwo@set@fontsize\set@fontsize
\def\set@fontsize#1#2#3{%
  \@defaultunits\labvtwo@size#2pt\relax\@nnil
  \advance\labvtwo@size by 1pt\relax
  \@defaultunits\labvtwo@skip#3pt\relax\@nnil
  \advance\labvtwo@skip by 1pt\relax
  \begingroup
    \edef\x{\endgroup
      \noexpand\labvtwo@set@fontsize{\unexpanded{#1}}{\strip@pt\labvtwo@size}{\the\labvtwo@skip}%
    }\x
}
\makeatother
\usetikzlibrary{fit,backgrounds,calc}

\makeatletter
\setlength{\beamer@frametopskip}{0mm}
\makeatother
\setbeamertemplate{frametitle}{%
  \vspace{4.2mm}%
  {\usebeamerfont{frametitle}\color{ink}\insertframetitle\par}%
  \vspace{0.4mm}%
  {\color{ruleblue}\rule{\textwidth}{0.7pt}\par}%
  \vspace{1.1mm}%
  \ifx\insertframesubtitle\@empty\else
    {\footnotesize\color{muted}\insertframesubtitle\par}%
    \vspace{1.2mm}%
  \fi
}

\definecolor{cviolet}{HTML}{F3E8F6}
\definecolor{cviolettt}{HTML}{7A458E}
\definecolor{cred}{HTML}{F8E6E6}
\definecolor{credtt}{HTML}{C44545}
\renewcommand{\usecardcolor}[1]{%
  \def\cardfill{cgreen}%
  \def\cardtitlec{cgreentt}%
  \ifstrequal{#1}{red}{%
    \def\cardfill{cred}%
    \def\cardtitlec{credtt}%
  }{}%
  \ifstrequal{#1}{orange}{%
    \def\cardfill{corange}%
    \def\cardtitlec{corangett}%
  }{}%
  \ifstrequal{#1}{blue}{%
    \def\cardfill{cblue}%
    \def\cardtitlec{cbluett}%
  }{}%
  \ifstrequal{#1}{green}{%
    \def\cardfill{cgreen}%
    \def\cardtitlec{cgreentt}%
  }{}%
  \ifstrequal{#1}{violet}{%
    \def\cardfill{cviolet}%
    \def\cardtitlec{cviolettt}%
  }{}%
  \ifstrequal{#1}{plain}{%
    \def\cardfill{cplain}%
    \def\cardtitlec{cplaintt}%
  }{}%
  \ifstrequal{#1}{gray}{%
    \def\cardfill{cplain}%
    \def\cardtitlec{cplaintt}%
  }{}%
}

\kicker{サーベイ}
\title{安全保障と好み選択を\\両立する局所プランナ}
\subtitle{CATNAV と SEA-Nav}
\date{2026年11月}

\newcommand{\li}{%
  \ifvmode\else\par\fi
  \noindent\hangindent=1.15em\hangafter=1
  \makebox[1.15em][l]{\textcolor{ruleblue}{\textbullet}}\ignorespaces}
\newcommand{\lii}{%
  \ifvmode\else\par\fi
  \noindent\hspace{1.0em}\hangindent=2.15em\hangafter=1
  \makebox[1.05em][l]{\textcolor{muted}{\textendash}}\ignorespaces}

\newcommand{\cardbalance}{\par\vspace*{-\topskip}\vspace{0pt plus 1fill}}
\newcommand{\takeawaybalance}[1]{%
  \par\vspace{0.15em plus 1fill}%
  \begingroup
  \setlength{\fboxsep}{5pt}%
  \noindent\colorbox{ink}{%
    \parbox[c]{\dimexpr\textwidth-10pt\relax}{%
      \color{white}\footnotesize #1}}%
  \endgroup
}

\RenewDocumentCommand{\infocard}{O{green} o m m}{%
  \usecardcolor{#1}%
  \IfValueTF{#2}{\def\cardminh{#2}}{\def\cardminh{0pt}}%
  \begin{tikzpicture}[baseline=(card.north)]
    \node[
      anchor=north west,
      fill=\cardfill,
      text width=\dimexpr\linewidth-14pt\relax,
      minimum height=\cardminh,
      inner xsep=6.5pt,
      inner ysep=4pt,
      align=flush left
    ] (card) {%
      {\color{\cardtitlec}\bfseries #3}\par\vspace{0.15em}%
      {\color{body}\footnotesize #4\par}%
    };
  \end{tikzpicture}%
}

\begin{document}

\begin{frame}[plain]
  \labtitlepage{}
\end{frame}

\begin{frame}{研究テーマ ｜ ロボットの制御系}
  \centering
  \begin{tikzpicture}[node distance=6.5mm]
    \node[kworange, text width=24mm, minimum height=8mm] (c) {コントローラ};
    \node[kwgreen, text width=24mm, minimum height=8mm, right=of c] (p) {局所プランナ};
    \node[kwblue, text width=24mm, minimum height=8mm, right=of p] (g) {歩行モデル};
    \node[kwpurple, text width=24mm, minimum height=8mm, right=of g] (m) {モータ};
    \draw[kwarr] (c) -- (p)
      node[midway, below=4mm, font=\scriptsize, text=muted] {目標座標};
    \draw[kwarr] (p) -- (g)
      node[midway, below=4mm, font=\scriptsize, text=muted] {目標速度};
    \draw[kwarr] (g) -- (m)
      node[midway, below=4mm, font=\scriptsize, text=muted] {トルク};
  \end{tikzpicture}

  \vspace{1.6em}
  \begin{columns}[T,onlytextwidth]
    \begin{column}{0.32\textwidth}
      \infocard[orange][22mm]{入力}{%
        \li 目標は座標 $(x,y,\theta)$
        \li 局所プランナが速度 $(v_x,v_y,\omega)$ に変換する}
    \end{column}
    \hfill
    \begin{column}{0.32\textwidth}
      \infocard[blue][22mm]{歩行}{%
        \li 歩行モデルが速度に追従する
        \li 出力は各関節トルク $(\tau_0,\ldots,\tau_{12})$}
    \end{column}
    \hfill
    \begin{column}{0.32\textwidth}
      \infocard[green][22mm]{今回の対象}{%
        \li 見るのは局所プランナ
        \li 速度指令に，安全と好みの両方を載せる}
    \end{column}
  \end{columns}

  \takeaway{目標座標は，局所プランナ・歩行モデルを経てモータへ届く．}
\end{frame}

\begin{frame}{研究テーマ ｜ 局所プランナ}
  \centering
  \begin{tikzpicture}[node distance=6.5mm]
    \node[kworange, text width=24mm, minimum height=8mm, opacity=0.25] (c) {コントローラ};
    \node[kwgreen, text width=24mm, minimum height=8mm, right=of c] (p) {局所プランナ};
    \node[kwblue, text width=24mm, minimum height=8mm, right=of p, opacity=0.25] (g) {歩行モデル};
    \node[kwpurple, text width=24mm, minimum height=8mm, right=of g, opacity=0.25] (m) {モータ};
    \draw[kwarr] (c) -- (p)
      node[midway, below=4mm, font=\scriptsize, text=muted] {目標座標};
    \draw[kwarr] (p) -- (g)
      node[midway, below=4mm, font=\scriptsize, text=muted] {目標速度};
    \draw[kwarr] (g) -- (m)
      node[midway, below=4mm, font=\scriptsize, text=muted] {トルク};
  \end{tikzpicture}

  \vspace{1.1em}
  \begin{columns}[T,onlytextwidth]
    \begin{column}{0.49\textwidth}
      \infocard[red][28mm]{安全保障}{%
        \li 障害物や人にぶつからない速度を出す
        \li 密集した環境でも，機敏に通り抜ける
        \li 今回は SEA-Nav で見る}
    \end{column}
    \hfill
    \begin{column}{0.49\textwidth}
      \infocard[green][28mm]{好み選択}{%
        \li 通れる道のうち，歩き方の好みを選ぶ
        \li 紙を踏まない，道の右側を歩く，など
        \li 今回は CATNAV で見る}
    \end{column}
  \end{columns}

  \takeaway{目標速度を，安全保障と好み選択の両方に対応させる．}
\end{frame}

\begin{frame}{今回の2本}
  \infocard[green][24mm]{好み ｜ CATNAV}{%
    \li Cached Vision-Language Traversability for Efficient Zero-Shot Robot Navigation
    \li Potnis et al.，arXiv，2026
    \li VLM で意味的なリスクを見積もり，複数経路から好みの経路を選ぶ}
  \vspace{0.45em}
  \infocard[blue][24mm]{安全 ｜ SEA-Nav}{%
    \li Efficient Policy Learning for Safe and Agile Quadruped Navigation in Cluttered Environments
    \li Chen et al.，arXiv，2026
    \li 微分可能な安全フィルタを方策の中に入れ，密集環境を短時間で学ぶ}

  \takeaway{1本目は好みを経路選択で表し，2本目は安全を速度指令の制約で表す．}
\end{frame}

% ---------- CATNAV ----------

\begin{frame}{CATNAV}
  \source{Potnis et al., arXiv, 2026}
  \cardbalance
  \infocard[green]{何をするか}{%
    \li RGB-D から，ロボットの体格に応じた走破コストマップをゼロショットで作る
    \li 似たシーンの評価はキャッシュし，VLM への問い合わせを減らす
    \li コストで引いた複数経路を画像に乗せ，VLM に安全な理由つきで選ばせる}
  \vspace{0.55em}
  \centering
  \begin{tikzpicture}[font=\footnotesize, box/.style={minimum height=11mm, text width=32mm, inner sep=2pt}]
    \node[kwblue, box] (a) {ゼロショット認識};
    \node[kworange, box, right=8mm of a] (b) {リスク数値の経路生成};
    \node[kwgreen, box, right=8mm of b] (c) {柔軟な経路選択};
    \draw[kwarr] (a) -- (b);
    \draw[kwarr] (b) -- (c);
  \end{tikzpicture}
  \takeawaybalance{認識，経路生成，経路選択の順で行動が決まる．}
\end{frame}

\begin{frame}{背景・課題}
  \source{Potnis et al., arXiv, 2026}
  \infocard[orange][20mm]{背景}{%
    \li 走破リスクは，ロボットの物理的な能力に対して評価する必要がある
    \li 屋外は動的障害物と地形，屋内は人と壊れやすいものへの配慮が要る}
  \vspace{0.35em}
  \begin{columns}[T,onlytextwidth]
    \begin{column}{0.32\textwidth}
      \infocard[plain][24mm]{学習ベース}{%
        \li 大量の学習データが要る
        \li 分布外に弱い}
    \end{column}
    \hfill
    \begin{column}{0.32\textwidth}
      \infocard[blue][24mm]{VLM}{%
        \li コスト化すると思考の柔軟さを捨てがち
        \li 毎回問うと計算が重い}
    \end{column}
    \hfill
    \begin{column}{0.32\textwidth}
      \infocard[orange][24mm]{VLA}{%
        \li 速度を直接出せる
        \li なぜその道かを人が読みにくい}
    \end{column}
  \end{columns}

  \takeaway{少ないデータと計算で，数値と人間らしい判断の両方を残したい．}
\end{frame}

\begin{frame}{課題の整理}
  \source{Potnis et al., arXiv, 2026}
  \cardbalance
  \infocard[blue][14mm]{1}{\li 少ないアノテーションでも，高い推論能力が欲しい}
  \vspace{0.3em}
  \infocard[orange][14mm]{2}{\li 少ない計算資源でも，高い周波数で動いてほしい}
  \vspace{0.3em}
  \infocard[green][14mm]{3}{\li 数値だけでなく，社会的で柔軟な判断をしてほしい}
  \takeawaybalance{できれば，その判断が人にとって解釈可能であってほしい．}
\end{frame}

\begin{frame}{貢献}
  \source{Potnis et al., arXiv, 2026}
  \cardbalance
  \infocard[blue][16mm]{貢献 1 ｜ ゼロショットのリスク}{%
    \li タスク固有の学習はせず，VLM の推論で意味的リスク値を推定する}
  \vspace{0.28em}
  \infocard[orange][16mm]{貢献 2 ｜ 評価の再利用}{%
    \li 画像特徴から類似シーンを見つけ，リスク評価をキャッシュして問い合わせを減らす}
  \vspace{0.28em}
  \infocard[green][16mm]{貢献 3 ｜ 経路の選択}{%
    \li RGB に乗せた複数経路から，視覚的に安全な経路を VLM に選ばせる}
  \takeawaybalance{VLM の柔軟な推論を残したまま，必要な問い合わせを減らす．}
\end{frame}

\begin{frame}{方法 ｜ 三つの段階}
  \source{Potnis et al., arXiv, 2026}
  \centering
  \begin{tikzpicture}[
      scale=0.82, transform shape,
      font=\scriptsize,
      box/.style={align=center, inner sep=2.2pt, minimum height=8.2mm}
    ]
    \node[kwblue, box, text width=28mm] (img) at (0,1.15) {RGB-D};
    \node[kwblue, box, text width=28mm] (clip) at (0,0) {CLIP 特徴量};
    \node[kwblue, box, text width=28mm] (nov) at (0,-1.15) {新規性チェック};
    \node[kworange, box, text width=30mm] (vlm) at (4.55,0.7) {VLM\\物体名とリスク};
    \node[kworange, box, text width=30mm] (cache) at (4.55,-0.85) {キャッシュ\\近傍 $k$ 件};
    \node[kwgreen, box, text width=28mm] (seg) at (8.7,0.7) {CLIP-Seg\\コストマップ};
    \node[kwgreen, box, text width=28mm] (path) at (8.7,-0.85) {T-RRT\\複数経路};
    \node[kwpurple, box, text width=26mm, minimum height=18mm] (sel) at (12.55,0) {VLM が\\経路を選ぶ};
    \draw[kwarr] (img) -- (clip);
    \draw[kwarr] (clip) -- (nov);
    \draw[kwarr] (nov.east) -- (vlm.west);
    \draw[kwarr] (nov.east) -- (cache.west);
    \draw[kwarr] (vlm) -- (cache);
    \draw[kwarr] (cache.east) -- (seg.west);
    \draw[kwarr] (seg) -- (path);
    \draw[kwarr] (path.east) -- (sel.south west);
    \begin{scope}[on background layer]
      \node[fill=cblue, rounded corners=1.5pt, inner sep=1.6mm, fit=(img)(nov)] (g2) {};
      \node[fill=corange, rounded corners=1.5pt, inner sep=1.6mm, fit=(vlm)(cache)] (g1) {};
      \node[fill=cgreen, rounded corners=1.5pt, inner sep=1.6mm, fit=(seg)(path)] (g3) {};
    \end{scope}
    \node[font=\scriptsize\bfseries, text=cbluett, anchor=north] at ([yshift=-1.6mm]g2.south) {貢献 2};
    \node[font=\scriptsize, text=muted, anchor=north] at ([yshift=-5.4mm]g2.south) {似ていれば聞かない};
    \node[font=\scriptsize\bfseries, text=corangett, anchor=north] at ([yshift=-1.6mm]g1.south) {貢献 1};
    \node[font=\scriptsize, text=muted, anchor=north] at ([yshift=-5.4mm]g1.south) {リスクは VLM};
    \node[font=\scriptsize\bfseries, text=cgreentt, anchor=north] at ([yshift=-1.6mm]g3.south) {経路生成};
    \node[font=\scriptsize\bfseries, text=cviolettt, anchor=north] at ([yshift=-1.6mm]g3.south -| sel.center) {貢献 3};
    \node[font=\scriptsize, text=muted, anchor=north] at ([yshift=-5.4mm]g3.south -| sel.center) {見た目で選ぶ};
  \end{tikzpicture}

  \takeaway{新規のシーンだけ VLM に聞き，似たシーンはキャッシュし，最後に経路を見る．}
\end{frame}

\begin{frame}{貢献 1 ｜ VLM によるリスク推定}
  \source{Potnis et al., arXiv, 2026}
  \begin{columns}[T,onlytextwidth]
    \begin{column}{0.56\textwidth}
      \infocard[orange][18mm]{聞く内容}{%
        \li 画像中の物体名と，そのリスクを出させる
        \li コストテーブルは $[\text{物体名},\,\text{リスク値}]$}
      \vspace{0.28em}
      \infocard[blue][18mm]{いつ聞くか}{%
        \li 処理は重く，毎回だと遅延が出る
        \li シーンが新しいときだけ問い合わせる}
    \end{column}
    \hfill
    \begin{column}{0.42\textwidth}
      \infocard[green][40mm]{ゼロショット}{%
        \li タスク用の学習データは使わない
        \li ロボットの体格に対するリスクを，VLM の知識で見積もる
        \li 人や構造物は高め，路面は低め，という差が付く}
    \end{column}
  \end{columns}

  \takeaway{VLM の知識だけでリスクを見積もれるが，毎回は処理したくない．}
\end{frame}

\begin{frame}{貢献 2 ｜ 画像特徴で問い合わせを控える}
  \source{Potnis et al., arXiv, 2026}
  \infocard[blue][16mm]{CLIP Encoder}{%
    \li 画像を 768 次元の特徴量にし，キャッシュ内のベクトルと比較する}
  \vspace{0.28em}
  \infocard[orange][18mm]{新規性の閾値}{%
    \li 閾値より特徴が違うときだけ，物体検出とリスク推定を実行する
    \li それ以外は，近傍 $k$ 件のリスク値を再利用する}
  \vspace{0.28em}
  \infocard[green][14mm]{効果}{%
    \li 似たシーンでは VLM に聞かない．聞いた結果はキャッシュへ追加する}

  \takeaway{画像特徴が近いときは，過去のリスク評価をそのまま使う．}
\end{frame}

\begin{frame}{CLIP-Seg ｜ コストマップ}
  \source{Potnis et al., arXiv, 2026}
  \begin{columns}[T,onlytextwidth]
    \begin{column}{0.49\textwidth}
      \infocard[blue][22mm]{セグメンテーション}{%
        \li コストテーブルの物体名との反応値を画素ごとに得る
        \li 反応が最大の物体がその画素だとみなす
        \li 小さいマップはバイリニア補間で画像サイズまで伸ばす}
    \end{column}
    \hfill
    \begin{column}{0.49\textwidth}
      \infocard[green][22mm]{2D マップ}{%
        \li 各画素のクラスに，対応するリスク値を乗せる
        \li 深度で 3D から 2D のコストマップへ落とす
        \li このマップの上で経路を引く}
    \end{column}
  \end{columns}

  \vspace{0.4em}
  \infocard[orange]{読み方}{%
    \li 舗装は低リスク，草は中程度，歩行者は高リスク，のように領域ごとのコストになる}

  \takeaway{物体名への反応の最大値で，ゼロショットに領域を分けてコストを乗せる．}
\end{frame}

\begin{frame}{軌道生成と軌道選択}
  \source{Potnis et al., arXiv, 2026}
  \begin{columns}[T,onlytextwidth]
    \begin{column}{0.48\textwidth}
      \infocard[plain][11mm]{1．中央}{\li 遷移ベース RRT（T-RRT）で経路を作る}
      \vspace{0.18em}
      \infocard[blue][11mm]{2．左}{\li 安全条件の範囲で，最大 $d$ だけ左へずらす}
      \vspace{0.18em}
      \infocard[blue][11mm]{3．右}{\li 同じ条件で，最大 $d$ だけ右へずらす}
    \end{column}
    \hfill
    \begin{column}{0.50\textwidth}
      \infocard[orange][18mm]{4．リスクトラック}{%
        \li 安全閾値を緩めた条件で，T-RRT をやり直す}
      \vspace{0.22em}
      \infocard[green][22mm]{VLM に選ばせる}{%
        \li 経路を RGB 画像へ描く
        \li 道の種類と，安全な理由を考えさせる
        \li 行動の指示（右側を歩く，紙を踏まない）もここで効く}
    \end{column}
  \end{columns}

  \takeaway{数値で候補を作り，見た目の理由で一本に絞る．}
\end{frame}

\begin{frame}{結果 ｜ コストマップ}
  \source{Potnis et al., arXiv, 2026}
  \cardbalance
  \begin{columns}[T,onlytextwidth]
    \begin{column}{0.32\textwidth}
      \infocard[green][28mm]{低い}{%
        \li 舗装などの歩きやすい地面
        \li 例としてリスク $0.1$}
    \end{column}
    \hfill
    \begin{column}{0.32\textwidth}
      \infocard[orange][28mm]{中間}{%
        \li 草など，一部は通れる領域
        \li 例としてリスク $0.3$}
    \end{column}
    \hfill
    \begin{column}{0.32\textwidth}
      \infocard[red][28mm]{高い}{%
        \li 人や構造物
        \li 例としてリスク $0.7$}
    \end{column}
  \end{columns}
  \takeawaybalance{セグメントされた領域ごとにコストが付き，人や構造物は高めになる．}
\end{frame}

\begin{frame}{結果 ｜ VLM への問い合わせ}
  \source{Potnis et al., arXiv, 2026, Table~II}
  \vspace{0.15em}
  \centering
  {\scriptsize
  \begin{tabular}{lccc}
    \toprule
    & Test 1 & Test 2 & Test 3 \\
    \midrule
    近傍数 $k$ & 1 & 5 & 5 \\
    新規性の閾値 $\gamma$ & 0.55 & 0.1 & 0.55 \\
    シーン問い合わせ & 16 & 125 & 9 \\
    経路選択の呼び出し & 156 & 295 & 157 \\
    VLM 問い合わせ合計 & 172 & 420 & 166 \\
    問合せ (Q/s) & 0.004 & 0.042 & 0.006 \\
    平均遅延 (s) & 3.937 & 17.35 & 3.717 \\
    \bottomrule
  \end{tabular}
  }

  \vspace{0.35em}
  \infocard[orange]{読み方}{%
    \li Test 2（$\gamma=0.1$）はほぼ全画像を新規とみなし，問い合わせが約 10 倍になる
    \li Test 3（$k=5$，$\gamma=0.55$）は Test 2 に比べ，VLM への問い合わせを 85.7\% 減らす}

  \takeaway{新規性の閾値を調整すると，画像の問い合わせ回数を減らせる．}
\end{frame}

\begin{frame}{結果 ｜ 柔軟な経路選択}
  \source{Potnis et al., arXiv, 2026, Table~III．各タスク 10 試行}
  \vspace{0.1em}
  \centering
  {\scriptsize
  \begin{tabular}{llcccc}
    \toprule
    & & 到達 (\%) & ゴール距離 & 衝突 (\%) & 行動違反 (\%) \\
    \midrule
    \multirow{2}{*}{道の右側} & OmniVLA & 60 & 2.7\,m & 20 & 40 \\
    & CATNAV & 70 & 2.8\,m & 10 & 30 \\
    \multirow{2}{*}{道の中央} & OmniVLA & 50 & 2.3\,m & -- & 20 \\
    & CATNAV & 70 & 2.25\,m & -- & 20 \\
    \multirow{2}{*}{人が横切る} & OmniVLA & 60 & 0.4\,m & 20 & -- \\
    & CATNAV & 70 & 0.3\,m & 20 & -- \\
    \multirow{2}{*}{ベンチ回避} & OmniVLA & 30 & 0.8\,m & 20 & -- \\
    & CATNAV & 40 & 0.7\,m & 30 & -- \\
    \multirow{2}{*}{床の紙} & OmniVLA & 90 & 0.1\,m & 0 & 70 \\
    & CATNAV & 90 & 0.1\,m & 0 & 30 \\
    \bottomrule
  \end{tabular}
  }

  \takeaway{到達率は平均 68\% 対 58\%．紙をよけるなど，行動違反は 33\% 少ない．}
\end{frame}

\begin{frame}{まとめ ｜ CATNAV}
  \source{Potnis et al., arXiv, 2026}
  \begin{columns}[T,onlytextwidth]
    \begin{column}{0.49\textwidth}
      \infocard[green][46mm]{貢献}{%
        \li タスク学習なしで，意味的リスクをゼロショット推定
        \li 類似シーンの評価を再利用し，問い合わせを 85.7\% 削減
        \li 画像上の複数経路から，理由つきで一本を選ぶ
        \li 好みは，経路選択の指示として入る}
    \end{column}
    \hfill
    \begin{column}{0.49\textwidth}
      \infocard[orange][46mm]{リミテーション}{%
        \li 2D--3D 変換でノイズが出る
        \li 軌道の固定パラメータが数か所ある
        \li CLIP-Seg など，基盤モデルの能力に依存する
        \li 細い障害物は遠方でリスクを過小評価しうる}
    \end{column}
  \end{columns}

  \takeaway{好みはコストの数値と，見た目による経路選択の二段で表現される．}
\end{frame}

% ---------- SEA-Nav ----------

\begin{frame}{SEA-Nav}
  \source{Chen et al., arXiv, 2026}
  \cardbalance
  \infocard[blue]{何をするか}{%
    \li 障害物が密集した環境で，四足の速度指令を強化学習する
    \li Actor の出力を，微分可能な CBF で安全な速度へ最小限ずらす
    \li Critic はずらした後の速度を評価するので，安全フィルタの介入も学習に入る}
  \vspace{0.4em}
  \centering
  \begin{tikzpicture}[font=\footnotesize, box/.style={minimum height=10mm, text width=28mm}]
    \node[kwblue, box] (a) {Actor};
    \node[kworange, box, right=7mm of a] (b) {LSE-CBF};
    \node[kwgreen, box, right=7mm of b] (c) {安全な速度};
    \node[kwpurple, box, right=7mm of c] (d) {Critic};
    \draw[kwarr] (a) -- (b);
    \draw[kwarr] (b) -- (c);
    \draw[kwarr] (c) -- (d);
  \end{tikzpicture}
  \takeawaybalance{安全範囲へ動かした行動を，方策自身が学習する．}
\end{frame}

\begin{frame}{背景}
  \source{Chen et al., arXiv, 2026}
  \infocard[orange][18mm]{自立ナビゲーション}{%
    \li 密な障害物の中の移動は難しく，安全性を上げる必要がある
    \li 大規模な事前学習で，局所環境の性能は上がってきた}
  \vspace{0.3em}
  \infocard[blue][20mm]{安全ナビと CBF の課題}{%
    \li Actor-Critic の外に CBF を置くと，方策は補正されたことを知らない
    \li 失敗でエピソードが終わると，その直前のデータを学習に使いづらい
    \li 取れる失敗サンプルの母数が少ない}

  \takeaway{フィルタの介入も含めて学び，衝突の前後を効率よくサンプルしたい．}
\end{frame}

\begin{frame}{従来手法}
  \source{Chen et al., arXiv, 2026}
  \begin{columns}[T,onlytextwidth]
    \begin{column}{0.32\textwidth}
      \infocard[plain][32mm]{模倣学習}{%
        \li 行動をまねる
        \li 未知だと失敗しやすい}
    \end{column}
    \hfill
    \begin{column}{0.32\textwidth}
      \infocard[orange][32mm]{深層強化学習}{%
        \li 機動性は出る
        \li 密集すると失敗が多い
        \li 安全制限を掛けにくい}
    \end{column}
    \hfill
    \begin{column}{0.32\textwidth}
      \infocard[blue][32mm]{後処理の安全フィルタ}{%
        \li CBF は Actor の外にある
        \li 方策は補正を知らない
        \li 介入を学習できない}
    \end{column}
  \end{columns}

  \takeaway{サンプル効率と，フィルタ込みの安全制御が同時に欲しい．}
\end{frame}

\begin{frame}{貢献}
  \source{Chen et al., arXiv, 2026}
  \cardbalance
  \infocard[blue][16mm]{貢献 1 ｜ End-to-End LSE-CBF}{%
    \li 安全条件を滑らかにし，学習時とテスト時の不一致をなくす}
  \vspace{0.28em}
  \infocard[orange][16mm]{貢献 2 ｜ 運動学的な正則化}{%
    \li 速度の範囲と変化の滑らかさを損失に入れ，実機へ安全に載せる}
  \vspace{0.28em}
  \infocard[green][16mm]{貢献 3 ｜ ACSI}{%
    \li 衝突前後を巻き戻して学び，サンプル効率のボトルネックを減らす}
  \takeawaybalance{安全層・正則化・衝突の再利用の三つで，短い学習でも走破する．}
\end{frame}

\begin{frame}{通常の Actor-Critic との違い}
  \source{Chen et al., arXiv, 2026}
  \centering
  \begin{tikzpicture}[x=9.2mm, y=9.2mm, font=\footnotesize, box/.style={align=center, minimum height=9mm, text width=22mm, inner sep=2pt}]
    \node[font=\scriptsize\bfseries, text=muted] at (0.15,2.05) {通常};
    \node[kw, draw=cplaintt, fill=cplain, box] (e1) at (0,0.85) {Environment};
    \node[kwblue, box] (a1) at (3.15,0.85) {Actor};
    \node[kwgreen, box] (c1) at (3.15,-0.55) {Critic};
    \draw[kwarr] (e1) -- (a1);
    \draw[kwarr] (a1.south) -- (c1.north)
      node[midway, right, font=\scriptsize, text=muted] {行動};

    \node[font=\scriptsize\bfseries, text=cbluett] at (7.55,2.05) {SEA-Nav};
    \node[kw, draw=cplaintt, fill=cplain, box] (e2) at (7.4,0.85) {Environment};
    \node[kwblue, box] (a2) at (10.45,0.85) {Actor};
    \node[kworange, box, text width=16mm] (s) at (13.15,0.85) {LSE-CBF};
    \node[kwgreen, box] (c2) at (10.45,-0.7) {Critic};
    \draw[kwarr] (e2) -- (a2);
    \draw[kwarr] (a2) -- (s);
    \draw[kwarr] (s.south) -- (s.south |- c2.east) -- (c2.east)
      node[pos=0.25, right, font=\scriptsize, text=muted] {安全な行動};
  \end{tikzpicture}

  \vspace{0.7em}
  \infocard[blue]{評価の対象}{%
    \li 通常は Actor の生の行動を Critic が見る
    \li SEA-Nav は，CBF で安全範囲へ動かした行動を Critic が評価する}

  \takeaway{安全フィルタは後処理ではなく，方策の一部として誤差が伝わる．}
\end{frame}

\begin{frame}{Actor の入出力}
  \source{Chen et al., arXiv, 2026}
  \begin{columns}[T,onlytextwidth]
    \begin{column}{0.49\textwidth}
      \infocard[blue][28mm]{入力}{%
        \li 行動履歴 10 フレームの特徴量
        \li 現在の観測
        \lii 41 本のレイで，障害物までの距離}
    \end{column}
    \hfill
    \begin{column}{0.49\textwidth}
      \infocard[green][28mm]{出力}{%
        \li 希望する速度指令 $\bar{u}$
        \li CBF の安全制約に関するパラメータ $\alpha$
        \lii 広い場所では大きく，狭い場所では小さく学ぶ}
    \end{column}
  \end{columns}

  \vspace{0.4em}
  \infocard[orange]{その後}{%
    \li LSE-CBF が $\bar{u}$ を安全な $u_s$ へ最小限動かし，Critic と損失は $u_s$ を見る}

  \takeaway{方策は速度だけでなく，どれだけ安全側に倒すか $\alpha$ も自分で決める．}
\end{frame}

\begin{frame}{貢献 1 ｜ LSE-CBF}
  \source{Chen et al., arXiv, 2026}
  {\footnotesize\color{body}
  各レイの安全条件は $h_i(x)=\rho_i-d_{\mathrm{safe}}$．$\rho_i$ は障害物までの距離．}

  \vspace{0.35em}
  \begin{columns}[T,onlytextwidth]
    \begin{column}{0.49\textwidth}
      \infocard[orange][24mm]{通常の CBF}{%
        \li $h(x)=\min_i h_i(x)$
        \li 最小のレイが入れ替わると不連続
        \li どの指令が悪いか，誤差を流しにくい}
    \end{column}
    \hfill
    \begin{column}{0.49\textwidth}
      \infocard[blue][24mm]{LSE-CBF}{%
        \li $\displaystyle h(x)=-\frac1k\ln\sum_i\exp(-k\,h_i(x))$
        \li 連続なので，悪い指令へ誤差が伝わる}
    \end{column}
  \end{columns}

  \vspace{0.3em}
  \infocard[green]{安全な速度}{%
    \li $\dot h(x,u)\ge -\alpha h(x)$ を満たす範囲で，$\|u_s-\bar{u}\|$ が最小の速度を出す
    \li Critic はその安全な速度を評価する}

  \takeaway{条件式を滑らかにすることで，学習時と実行時で同じ安全層を使える．}
\end{frame}

\begin{frame}{貢献 2 ｜ 損失}
  \source{Chen et al., arXiv, 2026}
  {\footnotesize
  \[
    L_{\mathrm{total}}=L_{\mathrm{PPO}}+\lambda_{\mathrm{shield}}L_{\mathrm{shield}}+\lambda_{\mathrm{reg}}L_{\mathrm{reg}}
  \]
  }
  \vspace{-0.6em}
  \begin{columns}[T,onlytextwidth]
    \begin{column}{0.33\textwidth}
      \infocard[blue][30mm]{$L_{\mathrm{PPO}}$}{%
        \li ゴール到達などの報酬
        \li 通常の方策更新}
    \end{column}
    \hfill
    \begin{column}{0.33\textwidth}
      \infocard[orange][30mm]{$L_{\mathrm{shield}}$}{%
        \li 安全層が元の指令を大きく変えないようにする
        \li $\alpha$ を小さくしすぎない}
    \end{column}
    \hfill
    \begin{column}{0.31\textwidth}
      \infocard[green][30mm]{$L_{\mathrm{reg}}$}{%
        \li 最大速度を超えない
        \li 速度と価値が急に変わらない}
    \end{column}
  \end{columns}

  \takeaway{正則化により，実機で安全かつ滑らかな速度指令を出せる．}
\end{frame}

\begin{frame}{貢献 3 ｜ ACSI}
  \source{Chen et al., arXiv, 2026}
  \infocard[blue][16mm]{バッファ}{%
    \li 直近 2 秒，100 ステップを保持し，衝突後には大きなペナルティを与える}
  \vspace{0.25em}
  \begin{columns}[T,onlytextwidth]
    \begin{column}{0.49\textwidth}
      \infocard[orange][22mm]{巻き戻し}{%
        \li $P_{\mathrm{reset}}$ で早期終了し，次のエピソードへ進む
        \li $P_{\mathrm{replay}}$ で 2 秒戻り，回避をやり直す}
    \end{column}
    \hfill
    \begin{column}{0.49\textwidth}
      \infocard[green][22mm]{リセット率}{%
        \li ゴール率に合わせて動的に変える
        \li $P_{\mathrm{reset}}=P_{\min}+(P_{\max}-P_{\min})\cdot\mathrm{clip}(L_{\mathrm{goal}}/1.5,\,0,\,1)$}
    \end{column}
  \end{columns}

  \takeaway{衝突した履歴を効率よく引くので，危険な局所を重点的に学べる．}
\end{frame}

\begin{frame}{結果 ｜ シミュレーション}
  \source{Chen et al., arXiv, 2026, Table~I．RTX 4090 で 2000 iter，約 1 時間}
  {\scriptsize\color{body} 10\,m 四方に障害物を置き，密度を Easy / Medium / Hard の 3 段階にする．数値は成功率 SR（\%）．}

  \vspace{0.35em}
  \centering
  {\scriptsize
  \begin{tabular}{lccc}
    \toprule
    & Easy & Medium & Hard \\
    \midrule
    SEA-Nav & \textbf{100.0} & \textbf{97.0} & \textbf{90.0} \\
    ACSI なし & 100.0 & 93.0 & 83.0 \\
    Shield なし & 97.7 & 90.0 & 74.3 \\
    $L_{\mathrm{reg}}$ なし & 96.3 & 82.3 & 57.0 \\
    \bottomrule
  \end{tabular}
  }

  \vspace{0.4em}
  \infocard[green]{Hard での衝突率}{%
    \li SEA-Nav は 5\%．ACSI なしは 8\%，Shield なしは 12\%，$L_{\mathrm{reg}}$ なしは 18\% まで上がる
    \li 三つの方法を入れると，難しい環境での走破がまとめて上がる}

  \takeaway{寄与は Easy では小さく，障害物が密な Hard で差が開く．}
\end{frame}

\begin{frame}{結果 ｜ 実機}
  \source{Chen et al., arXiv, 2026, Table~II．Unitree Go2}
  {\scriptsize\color{body}
  SEA-Nav は学習した歩行モデル，SEA-Nav-b は Unitree 公式の歩行モデル．
  SR は成功率（\%），AS は平均速度（m/s）．}

  \vspace{0.3em}
  \centering
  {\scriptsize
  \begin{tabular}{lcccccccc}
    \toprule
    & \multicolumn{2}{c}{密集した部屋} & \multicolumn{2}{c}{動的障害物} & \multicolumn{2}{c}{障害物コース} & \multicolumn{2}{c}{S 字} \\
    \cmidrule(lr){2-3}\cmidrule(lr){4-5}\cmidrule(lr){6-7}\cmidrule(lr){8-9}
    & SR & AS & SR & AS & SR & AS & SR & AS \\
    \midrule
    SEA-Nav & 100 & 1.6 & 90 & 1.5 & 100 & 1.2 & 100 & 1.4 \\
    SEA-Nav-b & 100 & 0.9 & 100 & 0.8 & 90 & 0.7 & 100 & 0.5 \\
    SLAM & 100 & 0.5 & 80 & 0.4 & 100 & 0.4 & 80 & 0.3 \\
    \bottomrule
  \end{tabular}
  }

  \vspace{0.35em}
  \infocard[blue]{読み方}{%
    \li 4 環境の多くで SEA-Nav の成功率が高い
    \li 公式歩行モデルや SLAM より速く，動く障害物でも SLAM より到達しやすい}

  \takeaway{短い学習の方策が，実機の密集環境でも速度と成功率を両立する．}
\end{frame}

\begin{frame}{まとめ ｜ SEA-Nav}
  \source{Chen et al., arXiv, 2026}
  \begin{columns}[T,onlytextwidth]
    \begin{column}{0.49\textwidth}
      \infocard[blue][46mm]{貢献}{%
        \li LSE-CBF を方策の中に通し，学習と実行の安全層を一致させる
        \li 運動学の正則化で，実機に出せる速度にする
        \li ACSI で衝突前後を再利用する
        \li Hard で成功率 90\%，約 1 時間の学習}
    \end{column}
    \hfill
    \begin{column}{0.49\textwidth}
      \infocard[orange][46mm]{リミテーション}{%
        \li 平らな路面を想定している
        \li 坂や段差は対象外
        \li 単純な局所解は抜けられる
        \li 複雑な行き止まりや迷路は未対応}
    \end{column}
  \end{columns}

  \takeaway{安全は，速度指令を通す微分可能な制約としてプランナに入る．}
\end{frame}

\begin{frame}{まとめ ｜ 2本のサーベイ}
  \cardbalance
  \begin{columns}[T,onlytextwidth]
    \begin{column}{0.49\textwidth}
      \infocard[green][50mm]{CATNAV ｜ 好み}{%
        \li VLM が意味的なリスクをゼロショットで付ける
        \li 似た画像はキャッシュし，問い合わせを減らす
        \li 複数経路を見て，指示に合う一本を選ぶ
        \li 歩き方の好みは，経路選択の理由として入る}
    \end{column}
    \hfill
    \begin{column}{0.49\textwidth}
      \infocard[blue][50mm]{SEA-Nav ｜ 安全}{%
        \li レイの距離から，安全な速度指令を直接学ぶ
        \li CBF は後処理ではなく，評価と損失に入る
        \li 衝突の巻き戻しで，密集環境を短時間で学ぶ
        \li 安全保障は，速度の制約として入る}
    \end{column}
  \end{columns}
  \takeawaybalance{好みは経路の選択，安全は速度の制約．局所プランナではこの二つを重ねたい．}
\end{frame}

\end{document}
